\documentclass{article}
\usepackage{graphicx} % Required for inserting images

\title{Interactive agent: control automation}
\author{Banele Mdluli}
\date{April 2026}

\begin{document}
\maketitle
\newpage
\section{Introduction}
The section contains low-level requirements for interactive and control summary agents. The interactive agent is used to answer complex queries from a user. The control summary agent will be used to complete an in-depth analysis of the control exceptions. 

\section{Low-level requirements}
\subsection{Interactive agents}
The architecture or workflow that will be used for the section is the evaluator-optimiser flow. The framework for developing the agents will be the OpenAI Agent SDK. The following are low-level requirements for the interactive agent.
\begin{enumerate}
    \item Create an agent that will serve as an assistant to a user who is querying about control performance information. 
    \item The assistant agent must consist of 3 agents. The following are the roles of each:
    \begin{enumerate}
        \item The request recipient agent should rewrite the query and add relevent context where applicable for the processing agent.
        \item The processing agent should analyse the request and call available API tools to compile an appropriate response.
        \item The evaluator agent should review if the response is adequate to share with user, and it has fulfilled the list of key requirements when providing a response. If the evaluator agent is satisfied with the response, the processing agent should return it to the user.
    \end{enumerate}
    \item The functional and non-functional requirements of the request recipient agent.
    \begin{enumerate}
        \item The agent should use RAG to collect relevant context for a query. The RAG method used must leverage LangChain libraries.
        \item The agent chunk performance should be subjected to precision@k, recall@k, and Mean reciprocal rank (MRR) thresholds.
        \item The agent should complete processing the request in less than 20 seconds.
    \end{enumerate}
    \item The functional and non-functional requirements of the processing agent.
    \begin{enumerate}
        \item If applicable, the agent must call APIs using the tool feature. To compile a response for the user.
        \item The agent's response to the user should be streamed.
        \item The agent's response should take less than 10 seconds.
        \item The agent is only allowed 2 attempts to create an adequate response; it should inform the user that it does not have relevant information to the question.
    \end{enumerate}
    \item The functional and non-functional requirements of the evaluator agent.
    \begin{enumerate}
        \item The agent should review the compiled feedback and determine if it is adequate to provide to the user.
        \item If the response is adequate, the reasons for approving the feedback should be recorded in the logs.
    \end{enumerate}

The token usage per agent for every query will be recorded to track overuse or underuse. User-agent interaction should be recorded for future reviews. 
\end{enumerate}

\subsection{Control summary agent}
The architecture or workflow that will be used for the section is the evaluator-optimiser flow. The framework for developing the agents will be the OpenAI Agent SDK. The following are low-level requirements for the control summary agent.

\begin{enumerate}
    \item Create an agent that will serve as an assistant to compile control performance summaries.
    \item The assistant agent must consist of 3 agents. The following are the roles of each:
    \begin{enumerate}
        \item The control data recipient agent should extract control exception data and retrieve relevant context(e.g., field name descriptions, control name,..) for the processing agent.
        \item The processing agent should analyse the request and call available API tools to compile an appropriate summary for the control performance.
        \item The evaluator agent should review if the summary is adequate to record on the control summary table, and it has fulfilled the list of key details when providing a summary. If the evaluator agent is satisfied with the summary, the processing agent should send the response to a summary table.
    \end{enumerate}
    \item The functional and non-functional requirements of the request recipient agent.
    \begin{enumerate}
        \item The agent should use RAG to collect relevant context for a query. The RAG method used must leverage LangChain libraries.
        \item The agent chunk performance should be subjected to precision@k, recall@k, and Mean reciprocal rank (MRR) thresholds.
        \item The agent should complete processing the request in less than 20 seconds.
    \end{enumerate}
    \item The functional and non-functional requirements of the processing agent.
    \begin{enumerate}
        \item If applicable, the agent must call APIs using the tool feature. To compile a summary for the control response.
        \item The agent's control summary response should have fewer than 500 characters.
        \item The agent's summary should take less than 3 minutes.
        \item The agent should use batch processing to analyse multiple control performance outcomes.
        \item The agent is only allowed 2 attempts to create an adequate summary; if it does not succeed, then it should send the reasons for inadequate responses.
    \end{enumerate}
    \item The functional and non-functional requirements of the evaluator agent.
    \begin{enumerate}
        \item The agent should review the compiled feedback and determine whether it is adequate to record on the control summary table or not.
        \item If the response is adequate, the reasons for approving the feedback should be recorded in the logs.
    \end{enumerate}
The token usage per agent for every query will be recorded to track overuse or underuse. User-agent interaction should be recorded for future reviews. 
\end{enumerate}



\end{document}
